\begin{solapartat}
\punts{0,25} A partir de la intensitat del camp gravitatori obtenim el mòdul de la força de la gravitatòria: $F_g = mg$

\punts{0,2} I la segona llei de Newton estableix que: $\vec{F} = m\vec{a}$

\punts{0,2} Com que sobre el satèl·lit només actua la força de la gravetat: $a = g$

\punts{0,2} D'altra banda, considerant que el satèl·lit descriu un moviment circular uniforme al voltant de la terra, la seva acceleració centrípeta és: $a = \dfrac{v^2}{R}$ o $a = \omega^2 R$

També s'ha de considerar correcte que directament s'expressi $mg = ma = m\dfrac{v^2}{R}$ o $mg = ma = m\omega^2 R$

\punts{0,2} Per tant $g = R\omega^2 = R\left(\dfrac{2\pi}{T}\right)^2 = \dfrac{4\pi^2}{T^2}R$

\punts{0,2} $g = \dfrac{4\pi^2}{T^2}R = \dfrac{4\pi^2}{(27,3\cdot 24\cdot 60\cdot 60)^2}384\times 10^{6} = 2,72\times 10^{-3}\ \mathrm{m/s^2}$
\end{solapartat}

\begin{solapartat}
\punts{0,2} Com que el camp gravitatori és conservatiu, l'energia mecànica a la superfície d'un planeta o d'un satèl·lit és igual a l'energia mecànica a l'infinit.

\punts{0,2} D'altra banda, l'energia mecànica mínima a l'infinit correspon a una energia cinètica nul·la, és a dir, l'energia mecànica mínima per poder escapar d'un objecte astronòmic és zero, $E_m(\infty) = 0$.

\punts{0,1} Per tant, atès que $E_m = E_P + E_C = 0 \Rightarrow E_C = -E_P$

Així l'energia cinètica mínima que cal subministrar és: $E_C(\text{superfície}) = -E_P(\text{superfície})$

\punts{0,1} Per tant, per a un objecte de massa $m$:
\[ E_C(\text{superfície}) = G\frac{M\cdot m}{R} \]
On $M$ i $R$ són la massa i el radi de l'objecte astronòmic, respectivament.

\punts{0,2} per la Terra \quad $E_C(\text{Terra}) = G\dfrac{M_T\cdot m}{R_T}$,

i per la Lluna \quad $E_C(\text{Lluna}) = G\dfrac{M_{Ll}\cdot m}{R_{Ll}}$.

\punts{0,45}
\[ \frac{E_C(\text{Terra})}{E_C(\text{Lluna})} = \frac{G\dfrac{M_T\cdot m}{R_T}}{G\dfrac{M_{Ll}\cdot m}{R_{Ll}}} = \frac{M_T}{M_{Ll}}\frac{R_{Ll}}{R_T} = \frac{81,3}{3,67} = 22,1 \]

Escapar de la Terra té, com a mínim, un cost energètic 22 vegades més gran que de la Lluna.
\end{solapartat}
